\documentclass[11pt]{article}
\usepackage[margin=2.4cm]{geometry}
\usepackage{fontspec}
\setmainfont{DejaVu Sans}
\usepackage{juliamono}
\usepackage{fvextra}
\usepackage{xcolor}
\usepackage{xurl}
\usepackage{hyperref}
\usepackage{amssymb}
\newcommand{\checkbox}{$\square$}
\newcommand{\answerchoices}{\checkbox\ YES \quad \checkbox\ NO \quad \checkbox\ CANNOT DETERMINE}
\newcommand{\decisionexplanation}{%
  \par\smallskip\textbf{If NO, give the reason for the mismatch (required). If
  CANNOT DETERMINE, state what context or evidence is missing:}\par
  \rule{\linewidth}{0.4pt}\par\vspace{7pt}\rule{\linewidth}{0.4pt}\par}
\setlength{\parindent}{0pt}\setlength{\parskip}{5pt}
\setlength{\emergencystretch}{2em}
\DefineVerbatimEnvironment{Vsrc}{Verbatim}{fontsize=\small,frame=single,
  breaklines=true,breakanywhere=true}
\DefineVerbatimEnvironment{Vlean}{Verbatim}{fontsize=\scriptsize,frame=single,
  numbers=left,numbersep=5pt,breaklines=true,breakanywhere=true}
\hypersetup{colorlinks=true,urlcolor=blue}

\begin{document}
\begin{center}{\LARGE Proof Verification Document}\\[2pt]{\large mse\_605342\_\_ans609345}\\[2pt]{\small math.StackExchange mined pool (confirmed)}\end{center}
\begin{center}\fcolorbox{blue!60!black}{blue!4}{\parbox{0.92\linewidth}{\textbf{Make three independent decisions.} Follow the evidence chain: original source $\longrightarrow$ challenged claim $\longrightarrow$ proposed counterexample $\longrightarrow$ Lean theorem. A record is valid only when all three links are correct. Complete the human sign-off before reading the automated review appendix; those earlier results are background, not an answer key.}}\end{center}
\textbf{Use only this document, except for the public source link.} Everything needed for Questions 2 and 3 is included here. For Question 1, open the public source link when one is shown. If there is no public link, the complete stored problem and solution are included here. If the information is still not enough, choose CANNOT DETERMINE.
For each question, choose \textbf{YES}, \textbf{NO}, or \textbf{CANNOT DETERMINE}. Choose CANNOT DETERMINE if the document does not provide enough information for you to decide; do not mark YES if you are uncertain. If you find a specific mismatch, choose NO and explain it in the space below that question. Do not rewrite the claim, add assumptions, or substitute a different counterexample.
\textbf{How copied text is shown.} The problem, quoted passage, and proposed counterexample come from the stored record. The Lean theorem statement and Lean files are copied exactly from the certificate. Some long lines continue on the next printed line; this is only page formatting and does not change the copied text. Line numbers are added only to help you find passages. If the stored record contains a character that cannot be printed, the document shows \texttt{[U+XXXX CONTROL]} in its place.
\begin{description}
\item[Challenged claim.] The exact statement from the source that this record says is false.
\item[Proposed counterexample.] A concrete case that should satisfy the claim's assumptions but make its conclusion false.
\item[Lean certificate.] A Lean file proving a formal theorem. Successful compilation proves the theorem as written, not that it is the right translation of the source.
\end{description}
\section*{Question 1 --- Does this proof verification document accurately represent the source?}
\textbf{What to decide.} Compare the challenged passage below with the original source---either the public link or the complete source text in this document. Answer YES only if they make the same claim and include the same assumptions or conditions. If this document changes the meaning, leaves out a condition, or points to a different passage, answer NO.
\subsection*{Original source}
\textbf{Source record:}
\begin{Vsrc}
qedbench_certificate_availability_triage/data/records_mined_pilot.jsonl
\end{Vsrc}
\textbf{Public source:}\par
\url{https://math.stackexchange.com/a/609345}\par
\subsection*{Problem or question from the source (verbatim)}
\begin{Vsrc}
Title: An algorithm for Gaussian integers

Question : Let $i^2=-1$. If a pair $(a+bi, c+di) (a,b,c,d\in\mathbb Z, a\not\equiv b, c\not\equiv d\ \text{(mod 2)})$ such that $a+bi$ and $c+di$ are relatively prime is given, then can we say that repeating the following operation leads $(1,1)$ in finitely many steps? 

  
  Suppose that $a+bi$ and $c+di$ are relatively prime if and only if $a+bi$ and $c+di$ don't have any common divisor except for $\pm 1,\pm i$.

  
  Operation : If a pair $(a+bi, c+di)$ is given (we may suppose that $N(a+bi)\le N(c+di)$ where $N(a+bi)=a^2+b^2$), the next pair is defined as $(a+bi, (c+di)+\epsilon (1+i)(a+bi))$ with an appropriate choice of $\epsilon=\pm 1,\pm i$.

Example : Noting that $3+2i$ and $4+i$ are relatively prime since each of them is a prime, we know that the pair $(3+2i,4+i)$ satisfies the condition above. Noting that $N(3+2i)\le N(4+i)$, 
$$\begin{align}
(3+2i,4+i) & \rightarrow (3+2i,-1+2i)\ \ \ \ (\because \ \ (4+i)+(+i)(1+i)(3+2i)=-1+2i)\\
 & \rightarrow (2-i, -1+2i)\ \ \ \  (\because \ \ (3+2i)+(+i)(1+i)(-1+2i)=2-i)\\ 
 & \rightarrow (2-i,-i)\ \ \ \  (\because \ \ (-1+2i)+(-i)(1+i)(2-i)=-i)\\
 & \rightarrow (1,-i)\ \ \ \  (\because \ \ (2-i)+(-1)(1+i)(-i)=1)\\ 
 & \rightarrow (1,1)\ \ \ \ (\because \ \ (-i)+(+1)(1+i)(+1)=1).
\end{align}$$

Motivation : I've been looking for a similar algorithm for Gaussian integers $\mathbb Z[i]$  as Euclidean Algorithm for rational integers $\mathbb Z$. Then, I reached the operation above. However, I'm facing difficulty for proving or disproving that repeating the operation leads $(1,1)$ for any pair. Can anyone help?

Edit 1 : As Gerry Myerson pointed out, the Euclidean algorithm works fine in the Gaussian integers. 

Edit 2 : I think if we can prove the following lemma, then we can say that repeating the operation leads $(1,1)$ in finitely many steps. However, I'm facing difficulty for proving the lemma. Can anyone help?

Lemma : If a pair $(a+bi, c+di) (a,b,c,d\in\mathbb Z, a\not\equiv b, c\not\equiv d\ \text{(mod 2)})$ such that $a+bi$ and $c+di$ are relatively prime with $N(a+bi)\le N(c+di)$ is given , then we can get 
$$N(c+di+\epsilon (1+i)(a+bi))\lt N(c+di)$$
with an appropriate choice of $\epsilon=\pm 1,\pm i$.
\end{Vsrc}
\subsection*{Challenged passage recorded from the source (stored text)}
\begin{Vsrc}
Lemma. If $a$ and $b$ are relatively prime, $N(a-\alpha b)\lt N(b)$.
\end{Vsrc}
\subsection*{Surrounding source lines (stored text with added line labels)}
The labels refer to lines in the solution text stored with the source record; the source-line content is not transliterated.
\begin{Vsrc}
0001 | We need some lemmas before concentrating on what matters. We shall prove the norm is always decreasing and is always nonzero in your situation. This would imply your claim if we can guarantee the algorithm never "chokes", i.e. we rule out the equality part. This can be made with additional informations.
0002 | 
0003 | Lemma. Let $t=1+i$. For every complex numbers $m,n$ it is possible to find a linear combination $\alpha$ of $\varepsilon t$ such that $$N(n-\alpha m)\leqslant N(m).$$
0004 | 
0005 | To prove the lemma, imagine the complex plane, the number $m$ and the number $n$, possibly very far away from $m$. Write the linear combinations of the numbers $mt, mit, -mt, -mit$. They cover the complex plane with squares of sides $|m|\sqrt{2}$. By the pigeonhole principle, $n$ is at most $|m|$ from another complex number, say, $\alpha m$. This implies $N(n-\alpha m)\leqslant N(m)$, as we wanted. 
0006 | 
0007 | Lemma. If $a$ and $b$ are relatively prime, $N(a-\alpha b)\lt N(b)$.
0008 | 
0009 | Suppose $N(a-\alpha b)=N(b)$. Following the idea of the first lemma, it follows that $a$ is in the center of a square of the covering we created. Now, consider the covering of all gaussian multiples of $b$. It contains all centers and all vertices, and, then, $a=kb$, contradiction.
0010 | 
0011 | Lemma. The norm $N(a-\alpha b)$ is always nonzero.
0012 | 
0013 | Suppose it is zero. Therefore, $a=\alpha b$, so they aren't relatively prime, contradiction.
0014 | 
0015 | Therefore, we proved the norm is decreasing and is always nonzero. This implies your algorithm always terminates in $(1,1)$.
\end{Vsrc}
\textbf{A YES means:} the proof verification document faithfully represents both the source context and the challenged claim, without dropping a condition or changing its meaning.
\section*{Question 2 --- Does the proposed counterexample disprove that claim?}
\textbf{Claim to test (repeated for convenience).}
\begin{Vsrc}
Lemma. If $a$ and $b$ are relatively prime, $N(a-\alpha b)\lt N(b)$.
\end{Vsrc}
\subsection*{How the record says the argument uses this claim}
This is the record's description, not an established fact; check it against the source.
\begin{Vsrc}
The proof's second lemma — "If $a$ and $b$ are relatively prime, $N(a-\alpha b)\lt N(b)$" — which supplies the strict norm decrease on which the final termination claim (line 0015) rests. The lemma is false whenever $b$ is a unit, because the proof's equality-elimination step ("$a=kb$, contradiction") is not a contradiction when $b$ is a unit: $b\mid a$ is fully compatible with $\gcd(a,b)$ being a unit. The proof never splits off this case.
\end{Vsrc}
\subsection*{Proposed counterexample (exact values or construction)}
\begin{Vsrc}
MINIMAL WITNESS (Gaussian integers Z[i], alpha = i):  a = -1,  b = -1.

 - b = -1 is a UNIT (N(b) = 1), so every common divisor of a and b divides -1 and is a unit: a and b ARE relatively prime. Explicit Bezout certificate: 0*a + (-1)*b = 1.
 - a - alpha*b = -1 - i*(-1) = -1 + i.
 - Recomputed quantities: N(a - alpha*b) = N(-1+i) = (-1)^2 + 1^2 = 2, while N(b) = N(-1) = 1.
 - The lemma asserts N(a - alpha*b) < N(b), i.e. 2 < 1. FALSE: 2 > 1. The purported strict norm decrease goes the wrong way, so the descent/termination argument at line 0015 has no basis in this case.

This is exactly the case flagged: b is a unit, so b | a is fully compatible with gcd(a,b) being a unit, and the proof's "a = k b, contradiction" elimination of equality does not apply.

The same failure occurs in every ring tested (b a unit, coefficients |.| <= 6):
 - Z[omega] (Eisenstein): a = -1, b = -1 -> N(-1+omega) = 3 >= N(b) = 1
 - Z[sqrt(-2)]: a = -2*sqrt(-2), b = -1 -> N(-sqrt(-2)) = 2 >= 1
 - Z[sqrt(-3)]: a = -2*sqrt(-3), b = -1 -> N = 3 >= 1
 - Z[sqrt(-5)]: a = -2 - sqrt(-5), b = -1 -> N(-2) = 4 >= 1
 - Z[sqrt(2)] (real, |N|): a = -2, b = -1 -> N(-2+sqrt2) = 2 >= 1
 - Z[phi]: a = phi, b = -1 -> N(2phi) = 4 >= 1
Phase-1 total: 7392 unit-b counterexamples across the seven rings.

The lemma also fails far beyond the unit case (Phase 2, 8172 counterexamples over certified-coprime pairs with neither element a unit), e.g. in Z[i]: a = -3-3i, b = -2+i, coprime via (-1+i)*a + (2+i)*b = 1, N(a - i*b) = N(-2-i) = 5 which is NOT < N(b) = 5; in Z[sqrt(-5)]: a = -sqrt(-5), b = -2, coprime via (-2+sqrt(-5))*a + (2+sqrt(-5))*b = 1, N(a - alpha*b) = N(sqrt(-5)) = 5 > N(b) = 4. Under the alternative reading (a, b rational integers, gcd = 1) the same unit case kills it: a = -40, b = -1, N(a - i*b) = 1601 >= N(b) = 1 (both the N(b) = b^2 and N(b) = |b| conventions).

STRUCTURAL REASON (setting-independent, so the witness does not depend on which alpha/ring the original proof intended): for any multiplicative norm with N(x) >= 1 for x != 0 and N(b) = 1 iff b is a unit, choosing b a unit makes a, b relatively prime for EVERY a, and the lemma would force N(a - alpha*b) < 1, i.e. a = alpha*b — false for every other a.
\end{Vsrc}
\subsection*{Why the record says the counterexample works}
Check this explanation independently; it is evidence to assess, not an answer key.
\begin{Vsrc}
The sentence asserts a universally quantified inequality over pairs of relatively prime ring elements: for all coprime a, b in Z[alpha], N(a - alpha*b) < N(b). A single coprime pair violating the inequality refutes it, so it is finitely refutable, and one was found by exhaustive exact-integer search.

Method and soundness safeguards:
1. Elements of R = Z[alpha] are integer pairs (x, y) meaning x + y*alpha, with alpha^2 = t*alpha - n; multiplication and norm N(x + y*alpha) = |x^2 + t*x*y + n*y^2| are exact integer operations. No floats, no randomness, stdlib only.
2. NEGATIVE CONTROL: before any search, the script verifies N(uv) = N(u)N(v) on all pairs with coefficients in [-3,3] in every ring (0 violations) and N(1) = 1; it aborts rather than report a witness if the arithmetic layer is inconsistent.
3. Coprimality is never assumed — it is CERTIFIED. Phase 1 uses b a unit with an explicit inverse found in the box (0*a + b^{-1}*b = 1). Phase 2 requires an explicit Bezout identity u*a + v*b = 1 found by exhaustive search; Bezout is a sound sufficient condition for relative primality in any commutative ring (any common divisor then divides 1). Pairs for which no certificate is found are SKIPPED, never claimed coprime, so no false witness can be produced. This also makes the result valid in Z[sqrt(-5)], which is not a UFD.
4. The reported witness is the non-degenerate one: a != 0 and a strict violation N(a - alpha*b) > N(b), not merely equality.

Because the blind protocol supplies only the sentence, the ring and alpha are not pinned down; the script therefore instantiates seven standard quadratic settings and finds the failure in all of them, and the unit-b failure mode is proved setting-independent by the norm-minimality argument above. So the refutation does not hinge on guessing the source's exact ring.

Runtime: ~8 seconds.
\end{Vsrc}
\subsection*{How to decide}
\begin{enumerate}
\item List every assumption that applies to the challenged claim.
\item Check that the proposed example satisfies every one of those assumptions.
\item Recompute the claimed outcome and the actual outcome using exact reasoning.
\item Confirm that the conflict is with the challenged claim itself, not an unrelated sentence.
\end{enumerate}
\textbf{A YES means:} the example is fully specified, satisfies all relevant assumptions, and makes the challenged claim false. If a necessary value, condition, or calculation is missing, choose CANNOT DETERMINE.
\section*{Question 3 --- Does the final Lean theorem statement verify this same counterexample?}
\textbf{What to decide.} A Lean file can compile while proving the wrong theorem. Compare these three links independently:
\begin{enumerate}
\item The Lean definitions represent the same mathematical objects, assumptions, and operations as the challenged claim.
\item The Lean construction uses the same concrete values or example shown in Question 2.
\item The final theorem statement asserts that the challenged claim fails on that example, not merely a weaker or nearby statement.
\end{enumerate}
\subsection*{Selected final Lean theorem statement (verbatim UTF-8 excerpt)}
This is an exact contiguous slice of the certificate file, from the theorem keyword through the original \texttt{:=} token where its proof begins. It is shown first so the formal target is not buried in the complete file.
\begin{Vlean}
theorem mse_605342__ans609345_refutation :
    bezoutCert a b bezoutU bezoutV = true
      ∧ mul alpha b = (0,-1)
      ∧ sub a (mul alpha b) = (-1,1)
      ∧ norm (sub a (mul alpha b)) = 2
      ∧ norm b = 1
      ∧ ¬ (norm (sub a (mul alpha b)) < norm b) :=
\end{Vlean}
\subsection*{Project-local imported Lean files (complete, verbatim UTF-8)}
This section follows the certificate's import commands recursively and prints every repository-owned Lean source file on which the final theorem statement depends. Each file is copied in full, including its definitions and controls. External Mathlib, Std, and Lean Init source is not duplicated.
No project-local imported Lean files were found. External Mathlib, Std, and Lean Init source is not duplicated in this proof verification document.
\subsection*{Complete counterexample Lean file (verbatim UTF-8)}
The complete Lean file is included below so you can check every definition and supporting result used by the final theorem statement. Use it together with the project-local imported Lean files above.
\textbf{Lean certificate file:}
\begin{Vsrc}
partB_feasibility/mined_certificates/mse_605342__ans609345.lean
\end{Vsrc}
\begin{Vlean}
/-!
# `mse_605342__ans609345` — coprimality does not force the claimed norm descent

Route: `SELF_CONTAINED_LEAN`; no V2 constructor models quadratic rings or their norms.
No imports are used.

We instantiate the sentence in the Gaussian integers with `alpha=i`, `a=-1`, `b=-1`.
The explicit Bézout identity `0*a + (-1)*b = 1` certifies coprimality.  Nevertheless
`a-alpha*b = -1+i`, whose norm is 2, while `N(b)=1`.  Thus the asserted strict decrease
would read `2<1`.

The negative control keeps the same unit `b` but takes `a=alpha*b`; then the difference
is zero and strict descent really holds.  This prevents a constant norm or broken
subtraction from vacuously certifying the counterexample.
-/

set_option maxRecDepth 10000
set_option maxHeartbeats 4000000

namespace MSE_605342_ans609345

/-- Gaussian integer `(re,im)`, denoting `re + im*i`. -/
abbrev GI := Int × Int

def add (z w : GI) : GI := (z.1 + w.1, z.2 + w.2)
def neg (z : GI) : GI := (-z.1, -z.2)
def sub (z w : GI) : GI := add z (neg w)
def mul (z w : GI) : GI := (z.1*w.1 - z.2*w.2, z.1*w.2 + z.2*w.1)
def norm (z : GI) : Int := z.1*z.1 + z.2*z.2

def zero : GI := (0,0)
def one : GI := (1,0)
def alpha : GI := (0,1)
def a : GI := (-1,0)
def b : GI := (-1,0)
def bezoutU : GI := zero
def bezoutV : GI := (-1,0)

def bezoutCert (x y u v : GI) : Bool := add (mul u x) (mul v y) == one

/-- The explicit coprime pair makes the alleged norm inequality false. -/
theorem mse_605342__ans609345_refutation :
    bezoutCert a b bezoutU bezoutV = true
      ∧ mul alpha b = (0,-1)
      ∧ sub a (mul alpha b) = (-1,1)
      ∧ norm (sub a (mul alpha b)) = 2
      ∧ norm b = 1
      ∧ ¬ (norm (sub a (mul alpha b)) < norm b) := by
  decide

/-! ### NEGATIVE CONTROL

At the unique value `a=alpha*b` the difference is zero and the descent succeeds.  The
same unit still supplies a Bézout certificate, so this is a non-vacuous positive instance
of exactly the modelled implication.
-/
def aGood : GI := mul alpha b

theorem mse_605342__ans609345_negative_control :
    bezoutCert aGood b bezoutU bezoutV = true
      ∧ sub aGood (mul alpha b) = zero
      ∧ norm (sub aGood (mul alpha b)) = 0
      ∧ norm (sub aGood (mul alpha b)) < norm b := by
  decide

end MSE_605342_ans609345


\end{Vlean}
\subsection*{Proof-safety information}
The automated scan found no unfinished proofs, new axiom declarations, unsafe declarations, external code hooks, or uses of Lean's compiled decision procedure. This scan concerns proof construction only; it does not establish that the theorem matches the source claim.
\textit{Technical scan counts:} native\_decide=0, sorry=0, axiom\_decl=0, unsafe=0, extern=0. \textit{Technical tier:} KERNEL-ONLY.
\textbf{A YES means:} Lean represents the same objects and assumptions, uses the same proposed counterexample, and the final theorem statement asserts the same contradiction. Successful compilation by itself is not enough.
\clearpage
\section*{Human sign-off}
Complete this block from your own review of Questions 1--3. Do not use the automated appendix below as an answer key.\par
\textbf{Reviewer name:} \rule{9cm}{0.4pt}\par
\textbf{Date:} \rule{4cm}{0.4pt}\par
\medskip\textbf{1. Does the proof verification document accurately represent the source and challenged claim?}\par
\answerchoices\par
\decisionexplanation
\medskip\textbf{2. Does the proposed counterexample satisfy the assumptions and make that claim false?}\par
\answerchoices\par
\decisionexplanation
\medskip\textbf{3. Does the final Lean theorem statement represent and verify that same counterexample?}\par
\answerchoices\par
\decisionexplanation
\medskip\textbf{Overall decision:}\par
\checkbox\ VALIDATED (all three answers are YES)\par
\checkbox\ REJECTED (at least one answer is NO)\par
\checkbox\ NEEDS ANOTHER REVIEWER (at least one CANNOT DETERMINE and no NO)\par
\medskip\textbf{Additional comments.} Use this space for any extra source, counterexample, calculation, or Lean details that support your decisions.\par\rule{\linewidth}{0.4pt}\par\vspace{7pt}\rule{\linewidth}{0.4pt}
\clearpage
\section*{Appendix --- Prior automated review (background only)}
\textbf{These results are not an answer key.} They come from the earlier Opus/Fable verification campaign and are printed only so an auditor can see what was previously checked. The human decisions above must be made independently. Raw status codes are retained in small type for traceability.
\textbf{Overall automated disposition:} Cleared by automated screening; an independent search also found a counterexample.
\par\textit{Audit code:} \nolinkurl{PASS_CORROBORATED}
\subsection*{Independent search with no access to the stored example}
\textbf{Result:} An independent search found a counterexample without seeing the stored example.
\par\textit{Audit code:} \nolinkurl{WITNESS_FOUND}
\par\textbf{Checker's explanation:}
\begin{Vsrc}
The quoted lemma asserts that a suitable multiplier alpha drives the norm strictly below N(b). Read with the alpha of this problem (alpha = eps(1+i), eps a unit -- the only multiplier the stated Operation permits), the claim is strictly stronger than the lemma the question actually needs (N(c+di+eps(1+i)(a+bi)) < N(c+di), i.e. '< N(a)' in the quote's letters), and it is false. Exact identity behind the failure: min over the four eps of N(a - eps(1+i)b) = N(a) + 2N(b) - 2(|x|+|y|) where x+iy = a*conj(b); (Shortened here; the complete note remains in the machine-check JSON.)
\end{Vsrc}
\par\textbf{Cases checked:}
\begin{Vsrc}
Deterministic stdlib-only Python 3, pure integer arithmetic (Gaussian integers as (re,im) int pairs; gcd by exact Gaussian Euclidean algorithm with integer round-to-nearest quotient floor((2p+q)/(2q)); no float influences any decision). (Shortened here; the complete note remains in the machine-check JSON.)
\end{Vsrc}
\subsection*{Independent re-check of the proposed counterexample}
\textbf{Result:} A separate checker recomputed the proposed counterexample and confirmed it.
\par\textit{Audit code:} \nolinkurl{CONFIRMED}
\subsection*{Independent review of the source-to-Lean correspondence}
\textbf{Result:} A separate reviewer judged that Lean proves the same failure using the same example.
\par\textit{Audit code:} \nolinkurl{MATCHES}
\par\textbf{Theorem summary from that reviewer:}
\begin{Vsrc}
In namespace MSE_605342_ans609345, Gaussian integers are modeled as GI := Int x Int with add/neg/sub/mul and norm z = z.1^2 + z.2^2, and constants alpha=(0,1)=i, a=(-1,0)=-1, b=(-1,0)=-1, bezoutU=(0,0), bezoutV=(-1,0). The main theorem mse_605342__ans609345_refutation asserts (by decide) the conjunction: bezoutCert a b bezoutU bezoutV = true (i.e. (Shortened here; the complete note remains in the machine-check JSON.)
\end{Vsrc}
\par\textbf{Difference noted by that reviewer:}
\begin{Vsrc}
No fidelity gap found; all three links hold. (1) Definitions: mul is the correct Gaussian product ((x1x2 - y1y2), (x1y2 + y1x2)), norm is the correct field norm x^2+y^2, sub is add-of-neg, alpha=(0,1)=i. These encode exactly the objects of the quoted lemma 'If a and b are relatively prime, N(a - alpha*b) < N(b)' as used in this Gaussian-integer problem, with neither weakening nor strengthening of the inequality (it is stated with strict <, and negated exactly). (Shortened here; the complete note remains in the machine-check JSON.)
\end{Vsrc}
\subsection*{Follow-up review}
No follow-up review was needed or recorded.
\textit{This appendix is only background. You do not need any separate machine-output files to answer the three questions.}
\end{document}